% Transcription — fonds Alexandre Grothendieck, shelfmark 162-5, batch 3 (pages 41-60).
% Established from the facsimile archives/batches/162-5/batch-03.pdf (a mirror of
% https://grothendieck.umontpellier.fr/162-5.pdf), per the skill
% transcribe-grothendieck. The batch file carries no cover sheet: PDF page N is
% archivists' page 40 + N, checked against the pencilled numbers bottom-left.
%
% What the batch holds, in the archivists' order.
%
%   41, 43-48  Seven leaves of the typescript « Tapis de Quillen », on which a hand
%              has worked in blue ink: the end of part II (page 41) and the whole
%              of part III, « Point de vue "motivique" en théorie du cobordisme »
%              (pages 43-48). The typed text is the same as that of pages 10 to
%              16 of the annotated photocopy that is folder 111; the ink layer
%              here is not identical to that folder's (folder 111 carries pencil
%              notes absent here — « NB H*(X) », « i.e. l'enveloppe Karoubienne »,
%              « = karoubienne + additive » — and this copy has ink the photocopy
%              lacks or shows less clearly, e.g. the label Γ = (p, id_X) on page
%              45 and the long marginal of page 48). The two were read
%              independently and each file says what its own leaves carry.
%              Typed pages carry no visible page number.
%   42, 49, 56 Blank. Skipped.
%   50-55      The offprint D. G. Quillen, « On the Representation of Hermitian
%              Forms as Sums of Squares », Inventiones math. 5 (1968), 237-242:
%              printed, not his, no mark of any hand on it. Skipped entirely,
%              without a note (it has no bearing on the typescript's mathematics).
%   57-60      A spirit-duplicated (blue) typescript in English by D. G. Quillen,
%              « On a conjecture of Adams », its pages 1 to 4, with Quillen's own
%              handwriting reproduced in the duplication (his name, a footnote on
%              page 58). Not transcribed. Page 58 alone appears, for one addition
%              in a darker ink foreign to the duplication (« dim E » as exponent),
%              anchored by a note; 57, 59, 60 carry nothing of the kind and are
%              absent. The typescript runs on past page 60 into batch 4 (pp. 61-68,
%              its pages 5-12 and bibliography), which carries no hand but
%              Quillen's duplicated one — the tick on p. 64 is in the
%              duplicator's blue — and is declared skipped in status.json.
%
% Attribution — as in folder 111, not asserted. The typed text is in the first
% person (« le rôle universel que j'ai envisagé plus haut », p. 47) and reports
% what « Quillen prouve »; the ink hand is not identified here either.
%
% Layers, following folder 111 so that the two copies of one text read alike.
% The typed text is set as running text; typed underlinings are \emph{} in prose
% and bold in mathematics ($\mathbf{B}$, $\mathbf{F}_{2}$, $\mathbf{Q}$,
% $\mathbf{S}$, $\underline{\mathrm{Hom}}$); typed interlinear insertions are
% typed text and are set in place without apparatus; typewriter x-outs are not
% reproduced. Everything in ink is apparatus: ink insertions into the typed line
% are \add{} (and \add{} is used for nothing else in this file), typed words
% struck by hand are \struck{}, and ink notes that stand apart from the line are
% \marginal{}. Symbols inked into blanks the machine left (arrows, ∈, ⊔, Λ,
% stars, tildes, overlines) are set as part of the formula without apparatus,
% with one note where it matters (page 45). Two blanks the machine left for Λ on
% page 46 were never inked and are marked \add{$\Lambda$}. The hand's small
% curled reference signs in the left margin (pages 41, 45, 46, 47) are not
% reproduced; the paragraphs they flag carry a note.
%
% Notation fixed for the batch: $F_{\cdot}$, $F^{\bullet}$ for the typed bold
% F. and F˙; $f_{\cdot}$, $f^{\bullet}$; $B_{\cdot}(X)$; $H^{\bullet}(X)$;
% $\Lambda = B_{\cdot}(\mathrm{pt})$; $\widetilde{\mathbf{B}}$ for the B the hand
% gives a tilde; $C^{\circ}$ is NOT supplied on page 44 where the typescript has
% a bare C.
%
% The typescript's slips are left as typed and not flagged one by one:
% « strcture », « ppouve », « thorème », « corespond », « interpétation »,
% « prédédente », « poihts », « rationelle », « néessairement », « Quillent »,
% « symmétrie », « flêches », « sss » (si et seulement si), « une B ».
%
% Crossing the batch boundaries. Page 41 opens mid-sentence (« qu'il conviendra
% de poser les définitions adéquates) ») — the sentence begins on page 40, in
% batch 2. Part III ends on page 48 with an ink paragraph closing on « ----- »;
% nothing of the typescript follows in this batch. Quillen's Adams typescript
% begun on page 57 continues into batch 4, which is skipped.
%
% Doubtful, not settled: the struck close of the ink paragraph on page 41 (the
% operation named before « des lacets »); the arrow labels of the composite on
% page 45; « entier ≥ 1 » in the ink insertion on Thom's generators, page 45;
% the struck word before $H^{1}(S^{1})$ and the exponent itself, page 47; the
% last two lines of the slanted marginal of page 48.
%
% Pass: Opus 5.5 (claude-opus-5-5), 2026-09-22 — first pass, unchecked against the pages by a human.
\documentclass[11pt,a4paper]{article}
% Le préambule commun à toutes les transcriptions du fonds Grothendieck.
%
% Il tient en une idée : **l'incertitude se compose, elle ne se résout pas.**
% Une transcription qui lisse un mot douteux a détruit la seule chose pour
% laquelle on la faisait — un lecteur doit pouvoir savoir, à chaque endroit,
% ce qui était sur le feuillet et ce qui a été ajouté. Les six macros qui
% suivent sont donc l'appareil critique, et non de la décoration.
%
% Les mêmes commandes sont reconnues par `scripts/render.mjs`, qui produit la
% vue de lecture du site. Une macro ajoutée ici doit l'être aussi là-bas,
% faute de quoi le rendu échouera bruyamment — ce qui est voulu : un rendu
% silencieusement faux est pire qu'un rendu absent.

\ProvidesPackage{grothendieck}[2026/08/08 Transcriptions du fonds Grothendieck]

\usepackage{amsmath,amssymb,amsthm}
\usepackage{mathrsfs}
\usepackage{stmaryrd}
% Les diagrammes commutatifs sont partout dans ce fonds : c'est la langue
% même de la géométrie algébrique de Grothendieck, et une transcription qui
% les rendrait en prose ne transcrirait rien.
\usepackage{tikz-cd}
% Français par défaut — c'est la langue des feuillets — mais l'anglais est
% chargé aussi : la traduction et le résumé sont des documents anglais, et la
% typographie française (espaces avant les deux-points) y serait une faute.
% Ces documents basculent par \selectlanguage{english} en tête de corps.
\usepackage[english,french]{babel}
\usepackage{fontspec}
\usepackage{geometry}
\geometry{a4paper,margin=25mm,marginparwidth=28mm}
\usepackage{marginnote}
\usepackage{xcolor}
% ulem, not soul. `soul`'s \st and \ul are built on a character-by-character
% scan that XeLaTeX's font machinery breaks: the document fails with an
% undefined control sequence rather than degrading. ulem does the same two jobs
% and is Unicode-engine safe. `normalem` stops it hijacking \emph, which the
% transcriptions use for Grothendieck's own emphasis.
\usepackage[normalem]{ulem}
\usepackage{hyperref}
\hypersetup{colorlinks=true,linkcolor=black,urlcolor=black,pdfborder={0 0 0}}

% --- Métadonnées du lot -----------------------------------------------------
% Reprises telles quelles par le script de rendu, qui les affiche en tête de
% la vue de lecture. Elles fixent ce que le fichier prétend couvrir : sans
% elles, une transcription flotte sans point d'attache dans un fonds de seize
% mille feuillets.

\newcommand{\folder}[1]{\gdef\grothFolder{#1}}
\newcommand{\batch}[1]{\gdef\grothBatch{#1}}
\newcommand{\leaves}[2]{\gdef\grothFirst{#1}\gdef\grothLast{#2}}
\newcommand{\foldertitle}[1]{\gdef\grothFoldertitle{#1}}
\gdef\grothFolder{?}\gdef\grothBatch{?}\gdef\grothFirst{?}\gdef\grothLast{?}\gdef\grothFoldertitle{}

% --- Le résumé --------------------------------------------------------------
%
% Il ouvre la lecture modernisée, sous le titre « Résumé », et il est composé
% en petit. Ce n'est pas un ornement : c'est le seul passage du document qui
% ne soit pas des mathématiques, et un lecteur doit pouvoir le reconnaître
% avant de l'avoir lu — puis le sauter s'il connaît déjà le sujet, ou s'y
% arrêter s'il ne le connaît pas. Le filet qui le suit marque la reprise du
% corps.

\newenvironment{resume}
  {\small\setlength{\parskip}{0.45em}}
  {\par\medskip\hrule\medskip}

% --- Mots-clés ---------------------------------------------------------------
%
% En anglais, en clôture du résumé de la lecture modernisée : le vocabulaire
% moderne sous lequel chercher ce que les feuillets construisent. Le manifeste
% du site (`npm run manifest`) les extrait pour étiqueter la cote entière —
% cette ligne est la source unique des tags, il n'y a pas de fichier à côté.
\newcommand{\keywords}[1]{\par\smallskip\noindent
  {\footnotesize\textsc{Keywords} --- \textit{#1}}\par}

% --- L'appareil critique ----------------------------------------------------

% Le numéro de feuillet, en marge. C'est l'ancre qui permet de régler un
% désaccord sur une formule en regardant *un* feuillet plutôt qu'une chemise
% de six cents. Le site s'en sert aussi pour tourner les pages du fac-similé
% quand on fait défiler la transcription.
\newcommand{\leaf}[1]{%
  \marginnote{\footnotesize\color{gray!70}\textsf{#1}}%
  \label{leaf:#1}\ignorespaces}

% Illisible : on ne devine pas. Un mot inventé qui se lit comme les autres est
% le pire résultat possible.
\newcommand{\ill}{\textcolor{red!60!black}{[\,\ldots\,]}}

% Lecture douteuse : le mot est donné, et signalé comme incertain.
\newcommand{\uncertain}[1]{\uline{#1}}

% Ajout de l'éditeur — une lettre manquante, un mot sous-entendu.
\newcommand{\add}[1]{\textcolor{blue!50!black}{[#1]}}

% Biffé par Grothendieck lui-même. Ce qu'il a rayé fait partie du document :
% une rature dit souvent où la pensée a changé de direction.
\newcommand{\struck}[1]{\sout{#1}}

% Note du transcripteur, sur l'état matériel du feuillet.
\newcommand{\note}[1]{\par\smallskip{\small\color{gray!60!black}\textit{[#1]}}\par\smallskip}

% Note marginale de Grothendieck, distincte du corps du texte.
\newcommand{\marginal}[1]{\par\smallskip\noindent\hspace{1em}%
  {\small\color{gray!60!black}#1}\par\smallskip}

% --- Titre ------------------------------------------------------------------

\AtBeginDocument{%
  \begin{center}
    {\large\bfseries Fonds Alexandre Grothendieck --- cote n\textsuperscript{o}~\grothFolder}\\[2pt]
    {\itshape\grothFoldertitle}\\[4pt]
    {\small feuillets \grothFirst--\grothLast\ (lot \grothBatch)}\\[2pt]
    {\footnotesize\color{gray!60!black}Transcription établie d'après le fac-similé numérisé
     par l'Université de Montpellier.\\
     Les lectures incertaines sont soulignées, les illisibles notées [\,\ldots\,],
     les ajouts entre crochets.}
  \end{center}
  \vspace{1em}}

\endinput


\folder{162-5}
\batch{3}
\pages{41}{60}
\dating{1968-[à partir de 1970] — le groupe « Dossiers rassemblés par Grothendieck sur différentes thématiques » (161-1 à 162-6) est daté [après 1961-vers 1977]}
\watermark{Édition de démonstration}
\foldertitle{Tapis de Quillen : tapuscrit et copies de tapuscrits (1968, s.d.), notes manuscrites (1968), tiré à part (1968)}

\begin{document}

\page{41}

\note{feuillets tapés, repris à l'encre. Le texte tapé est donné en clair ;
les insertions manuscrites dans la ligne sont données entre crochets, les mots
tapés biffés à la main sont barrés, les notes manuscrites détachées du texte
sont données à part, sans que cette présentation attribue la main. Le texte
tapé des pages 41 et 43 à 48 est celui des pages 10 à 16 de la copie
photographique du dossier 111 ; les interventions manuscrites ne sont pas
toutes les mêmes}

qu'il conviendra de poser les définitions adéquates). --- Peut-être Quillen
a-t-il une définition semi-simpliciale plus directe, sans passer par des
$n$-catégories, des $K_{i}$ et plus précisément du H-espace qui devrait
correspondre à $C$.

\marginal{oui}\note{dans la marge gauche, en regard de la phrase qui précède,
que la main marque d'un double trait vertical}

Ce n'est certainement pas en général l'ensemble semi-simplicial $SC$ associé à
$C$, puisque dans le cas d'une catégorie additive, $C$ est homotopiquement
trivial (en effet le foncteur identique est « homotope » au foncteur nul,
puisqu'il suffit d'avoir un homomorphisme d'un foncteur dans l'autre ; d'où une
équivalence d'homotopie entre $C$ et la catégorie nulle). Pourtant, il semble
bien que dans le cas où $C$ est une Gr-catégorie i.e.\ admet des inverses,
alors ce soit bien la construction « naïve » $SC$ qui est la bonne, puisqu'on
trouve bien un $\pi_{0}$ et un $\pi_{1}$ qui sont aussi les $\pi_{0}$ et
$\pi_{1}$ au sens catégorie sans Gr-structure.

\marginal{A expliquer : pour des Gr-$n$-catégories, les $\pi_{i}$ au sens
semi-simplicial sont aussi les $\pi_{i}$ au sens Gr-$n$-catégories, savoir les
objets à isom.\ près dans divers $\underline{\mathrm{Hom}}(1, 1)$ …\
\struck{Il y a aussi une interprétation de l'opération \uncertain{$\Omega$ :
\ill{} des lacets} en termes du foncteur d'opération du shift géométrique
évidente sur les $n$-catégories …}}\note{paragraphe à l'encre sous la frappe,
appelé, comme la phrase « i.e.\ admet des inverses … » qui le précède, par un
signe en marge et un trait vertical ; sa seconde moitié est barrée de traits
obliques}

\page{43}

\subsection*{III \quad Point de vue « motivique » en théorie du cobordisme.}

Soit $C$ la catégorie des variétés différentiables (pas nécessairement
orientables), les morphismes étant les classes d'homotopie d'applications
continues. Si $B$ est une catégorie, on s'intéresse aux couples $(F_{\cdot},
F^{\bullet})$ d'un foncteur covariant et d'un foncteur contravariant de $C$ dans
$B$, satisfaisant les conditions que pour tout $X \in \mathrm{Ob}\, C$, on a
$F_{\cdot}(X) = F^{\bullet}(X)$, et que si on a un produit fibré ordinaire

\begin{tikzcd}
 & X' \arrow[dl, "g'"'] \arrow[dr, "f'"] & \\
X \arrow[dr, "f"'] & & Y' \arrow[dl, "g"] \\
 & Y &
\end{tikzcd}

avec $f$ et $g$ transversaux, alors on a $g^{\bullet} f_{\cdot} = f'_{\cdot}
g'^{\bullet}$ (on pose $f_{\cdot} = F_{\cdot}(f)$, $f^{\bullet} =
F^{\bullet}(f)$). La question est de trouver une $B$ et un couple $(F_{\cdot},
F^{\bullet})$ qui soient « universels » dans un sens évident. Quillen prouve que
la construction suivante donne une telle solution.\note{le carré est tapé en
losange, $X'$ en haut et $Y$ en bas ; « les morphismes étant les classes
d'homotopie d'applications continues » et « (on pose …) » sont tapés en
interligne et mis en place par un trait à l'encre}

Pour toute variété $X \in \mathrm{ob}\, C$, on définit un anneau commutatif
$B(X)$ (l'anneau de bordisme de $X$) en prenant comme éléments les classes
d'éléments de $C$ au dessus de $X$, pour la relation de cobordisme : $f' : X'
\to X$ et $f'' : X'' \to X$ sont équivalents, s'il existe une variété à bord
$Z$, et un $f : Z \to X$, tel que $\mathrm{bord}\, Z = X' \sqcup X''$ et $f'$
et $f''$ sont induits par $f$. La somme de $B(X)$ est la somme disjointe (c'est
une loi de groupe, où tout élément est d'ordre \add{2}), le produit est donné
par le produit fibré, en utilisant le théorème de transversalité de Thom pour
bouger un peu $f''$, de sorte à le rendre transversal à $f'$.\note{le chiffre
de « d'ordre 2 » est récrit à l'encre sur une frappe biffée qui ne se lit pas}

\marginal{\emph{NB} L'anneau est de plus \emph{gradué} par la dimension
relative …}\note{à l'encre, en interligne}

Si on a $f : X \to Y$, on définit $f_{\cdot} : B(X) \to B(Y)$ de façon
évidente (toute variété sur $X$ est sur $Y$), et $f^{\bullet} : B(Y) \to B(X)$
par produit fibré (utilisant encore le th.\ de transversalité). De cette façon,
$B(X)$ devient un foncteur-anneau \add{(gradué)} contravariant en $X$, et un
foncteur groupe covariant en

\page{44}

$X$. De plus, on a la formule \add{(du type $g^{\bullet} f_{\cdot} =
f'_{\cdot} g'^{\bullet}$ ci-dessus, impliquant la formule)} de projection
habituelle. Ceci posé, les constructions bien connues en théorie des
correspondances permettent de définir une catégorie $\mathbf{B}$ dont les objets
sont ceux de $C$, et les Hom étant donnés par
\[
\mathrm{Hom}_{\mathbf{B}}(X, Y) = B(X \times Y) .
\]
On a un isomorphisme canonique $\mathbf{B} \simeq \mathbf{B}^{\circ}$, qui est
l'identité sur les objets et la symmétrie $B(X \times Y) \simeq B(Y \times X)$
sur les flêches, et un foncteur évident $F_{\cdot} : C \to \mathbf{B}$, qui est
l'identité sur les objets et donné par les graphes sur les flêches. Composé avec
la symmétrie précédente, on trouve un foncteur $F^{\bullet} : C \to
\mathbf{B}$.\note{la frappe porte ici $C$, sans le $^{\circ}$ qu'appelle un
foncteur contravariant ; rien n'a été ajouté à l'encre} On vérifie enfin que le
couple $(F_{\cdot}, F^{\bullet})$ satisfait aux conditions envisagées dans
l'alinéa 1 (ceci aussi devrait être formalisé au cas des correspondances
générales …). On vérifie enfin qu'il a la propriété universelle qu'il faut.
Lorsqu'on se donne un couple $(F'_{\cdot}, F'^{\bullet})$ à valeurs dans une
catégorie $B'$, il faut trouver une factorisation par $B$. Elle est toute
trouvée sur les objets, il faut la définir sur les flêches, donc définir, pour
$X$, $Y$ dans $C$, une application
\[
B(X \times Y) \longrightarrow \mathrm{Hom}(X', Y') . \tag{$*$}
\]
Mais si $Z$ est au dessus de $X \times Y$, d'où $p : Z \to X$ et $q : Z \to Y$,
on lui associe un élément du deuxième membre de $(*)$ comme le composé
\[
X' \xrightarrow{F'^{\bullet}(p)} Z' \xrightarrow{F'_{\cdot}(q)} Y' .
\]
Il faut prouver que l'élément obtenu en remplaçant $Z$ \add{$(= Z_{0})$} par un
$Z_{1}$ cobordant au dessus de $X$ \add{est le même}. Or si $g : T \to X \times
Y$ réalise le cobordisme, on voit (regardant seulement $T$ et la décomposition
$\mathrm{bord}(T) = Z_{0} \sqcup Z_{1}$) qu'il existe une variété
$\overline{T}$ sur $S^{1}$ (obtenue en « doublant » $T$) transversale au dessus
d'un couple de points $s_{0}$, $s_{1}$ de $S^{1}$, tels que $Z_{i}$ soit l'image
inverse \add{dans $\overline{T}$} de $s_{i}$ ($i = 0, 1$) ; d'autre part,

\page{45}

$g$ manifestement se prolonge en $\overline{g} : \overline{T} \to X \times Y$.
Cela nous ramène à prouver que le composé
\[
\overline{T}' \xrightarrow{\alpha_{i}} Z'_{i} \xrightarrow{\alpha_{i*}}
\overline{T}'
\]
ne dépend pas de $i = 0, 1$.\note{les deux étiquettes sont à l'encre ; la
première se lit $\alpha_{i}$, sans astérisque visible. Devant $\overline{T}'$ et
$Z'_{i}$ la frappe porte de petits signes corrigés, peut-être des $F'$, qu'on
ne reproduit pas} Ecrivons $X$ au lieu de $\overline{T}'$, et considérons le
diagramme cartésien d'inclusions

\begin{tikzcd}
X_{s} \arrow[r, "\alpha_{s}"] \arrow[d, "\alpha_{s}"'] & X \arrow[d, "{\Gamma = (p, \mathrm{id}_{X})}"] \\
X \arrow[r, "i_{s}"'] & S^{1} \times X
\end{tikzcd}

on a alors $\alpha_{s}^{*} \alpha_{s*} = \Gamma^{*} i_{s*}$, qui ne dépend pas
de $s$ puisque $i_{s}$ n'en dépend pas à homotopie près. --- Le reste est sans
doute l'AQT.\note{les flèches du carré, les $\alpha$, les astérisques et
l'étiquette $\Gamma = (p, \mathrm{id}_{X})$ sont portés à l'encre dans les
blancs de la frappe}

Sur la strcture des $B_{\cdot}(X)$. Ce sont tous des algèbres graduées sur
$B_{\cdot}(\mathrm{pt}) = \Lambda$, qui est une algèbre graduée à degrés
positifs, réduite à $\mathbf{F}_{2}$ en degré 0. C'est d'ailleurs une algèbre
de polynômes \add{(à une infinité de générateurs $x_{i}$ ($i$
\uncertain{entier $\geq 1$}, $\neq$ de la forme $2^{\nu} - 1$))} connue
…\note{« d'anneaux gradués » est tapé en interligne sous la ligne, sans place
marquée ; l'insertion à l'encre est appelée par un signe en marge}

On a un homomorphisme évident $B_{\cdot}(X) \to H^{\bullet}(X)$, nul en degrés
$> 0$, qui à tout $Z$ sur $X$ associe l'image de $1 \in H^{0}(Z)$ par
l'homomorphisme de Gysin associé à $Z \to X$. Cet homomorphisme par Thom est
surjectif (toute classe de cohomologie provient d'une image de variété compacte
…) ;

\marginal{en fait $H^{\bullet}(X) \simeq B_{\cdot}(X) / \Lambda^{+}
B_{\cdot}(X)$}\note{à l'encre, dans la marge gauche}

On ppouve même qu'il existe un relèvement fonctoriel en $X$ $H^{\bullet}(X)
\to B_{\cdot}(X)$, \add{(il est douteux qu'il soit)} compatible avec
multiplication, tel que l'on ait :\note{la main entoure « compatible avec
multiplication » d'une parenthèse et écrit au-dessus « il est douteux qu'il
soit » ; un point d'interrogation dans la marge}

\struck{a)} L'homomorphisme correspondant $H^{\bullet}(X)
\otimes_{\mathbf{F}_{2}} \Lambda \simeq B_{\cdot}(X)$ est un isomorphisme (ce
sera vrai pour tout relèvement individuellement …), \struck{et}

\struck{b)} [\add{Quillen ignore si} \struck{Le} relèvement est compatible
aussi avec les hom.\ de Gysin.]\note{les lettres a) et b) et le « et » sont
biffés à la main ; la seconde clause, mise entre crochets à l'encre, cesse
d'être une condition pour devenir une question. Un grand trait en arc, ou un
point d'interrogation, dans la marge en regard}

Ainsi, la catégorie $\mathbf{B}$ admet un foncteur pleinement fidèle dans la
catégorie des modules gradués libres de type fini sur $\Lambda$, les modules de
type fini qu'on obtient ainsi ayant des générateurs à degrés $< 0$ et
satisfaisant (pour $X$ connexe i.e.\ quand il y a un seul générateur de degré
zéro) la dualité de Poincaré. Quand on prend la catégorie quasi-abélienne
\add{$\widetilde{\mathbf{B}}$} associée à $\mathbf{B}$, on trouve donc tous les
modules gradués libres de type fini à générateurs en degrés $\leq 0$ seulement.
C'est une

\page{46}

catégorie additive (et même quasi-abélienne) avec $\otimes$ qui a toutes les
vertus, sauf d'être abélienne, et sauf d'admettre des \emph{Hom} internes. Pour
trouver ces derniers, on peut si on veut paraphraser la construction de la
catégorie des motifs, en considérant l'objet \add{$\Lambda(-1)$} de
$\widetilde{\mathbf{B}}$ conoyau de $p^{*} : e \to S$, où $p : S^{1} \to e$ est
l'application du cercle sur un point. C'est un objet de poids 1 (le poids est
défini par le degré des générateurs, changés de signe). On constate aussitôt que
le foncteur $M \mapsto M(-1) = M \otimes \Lambda(-1)$ de
$\widetilde{\mathbf{B}}$ dans lui-même est pleinement fidèle, ce qui permet
alors, par une construction formelle standard, de compléter
$\widetilde{\mathbf{B}}$ en une catégorie où $\Lambda(-1)$ devient inversible.
Il se trouve alors que dans cette dernière les $\underline{\mathrm{Hom}}(X, Y)$
internes existent et sont isomorphes à $\check{X} \otimes Y$ … en fait, on
trouve maintenant la catégorie de \emph{tous} les modules gradués libres de type
fini sur \add{$\Lambda$}. Bien sûr, cette catégorie n'est pas encore abélienne
pour autant, puisque \add{$\Lambda$} n'est pas un corps ; si on veut définir la
catégorie abélienne engendrée, on s'attend à trouver la catégorie des modules
de \struck{type} \add{présentation} finie sur $\Lambda$. (\struck{est-il
noethérien, i.e.\ le nombre de générateurs est-il fini ?}) \add{qui est
abélienne bien que $\Lambda$ ne soit pas noethérien, car $\Lambda$ est
\uncertain{$\varinjlim$} d'anneaux noeth.\ à homs de transition
\emph{plats})}.\note{les tildes sur $\mathbf{B}$, les $\Lambda$ et l'astérisque
de $p^{*}$ sont à l'encre ; la frappe écrit $S$ pour le cercle dans « $p^{*} : e
\to S$ ». Les deux \add{$\Lambda$} de « sur … » et « puisque … n'est pas un
corps » comblent des blancs que la frappe laissait et que l'encre n'a pas
remplis. « lim », à l'encre, est souligné d'un trait qui peut être la flèche
d'une limite inductive}

Cette construction ressemblerait à celle d'une théorie des motifs où on
n'aurait pas passé à l'équivalence homologique pour les cycles, mais où on
aurait gardé l'équivalence rationelle, --- qui est bien trop fine à beaucoup
d'égards. Pour trouver un équivalent plus étroit, il convient de reprendre la
théorie prédédente, on pourrait formuler le même Pb universel que plus haut,
mais en prenant l'équivalence homologique (via le graphe) des morphismes de
variétés, au lieu de l'équivalence homotopique. \struck{Mais alors il semble que
l'on trouve comme catégorie solution celle dont les objets sont les variétés,}

\page{47}

\struck{les morphismes \emph{toutes} les correspondances cohomologiques (en vertu
du thorème de Thom rappelé plus haut). On trouve donc un foncteur pleinement
fidèle de cette catégorie dans celle des vectoriels gradués positivement de dim.\
finie sur $\mathbf{F}_{2}$, qui est déjà abélienne, et si on veut rendre comme
plus haut $\mathbf{F}_{2}(-1)$ inversible, on trouve tous les vectoriels gradués
de dim.\ finie sur $\mathbf{F}_{2}$. C'est un peu trop trivial comme structure,
et on peut faire mieux,}\note{la dernière ligne de la page 46 est barrée d'un
trait, et les sept premières lignes de la page 47 de grands traits en zigzag}

Pour trouver une catégorie universelle, on associe à toute $X$
l'ensemble $S(X)$ des sous-variétés de $X$, modulo équivalence
homologique. Donc $S(X) \subset H^{\bullet}(X)$, et on prouve qu'on trouve en
fait un sous-anneau de $H^{\bullet}(X)$ (grâce à un théorème de Thom qui
caractérise $S(X)$ en termes d'opérations de Steenrod, cf plus bas). On prouve
de même (par le même théorème) que les $S(X \times Y)$ sont stables par
composition des correspondances cohomologiques (cela ne semble pas trivial,
puisque les $S(X)$ ne sont pas stables par Gysin …). D'où une catégorie additive
dont les objets sont ceux de $C$, les flêches sont les $S(X \times Y)$. Il
semble qu'elle doive jouer le rôle universel que j'ai envisagé plus haut. En
tous cas, sa définition géométrique est fort simple ! Nous allons en donner une
interpétation algébrique. Pour ceci, considérons l'anneau des opérations
cohomologiques « stables » \struck{i.e.\ des endomorphismes du foncteur $X
\mapsto H^{\bullet}(X)$ qui sont additifs, et qui en degré zéro se réduisent à
un multiple de l'identité.}\note{ce passage est barré de traits tapés, repassés
à l'encre d'un trait ondulé ; un signe dans la marge, et un trait vertical le
long du paragraphe} C'est un anneau gradué, avec application diagonale
anticommutative, i.e.\ commutative (car.\ 2 !), associative et unitaire. Son
dual gradué est une algèbre graduée anticommutative donc commutative, avec
application diagonale associative et unitaire ; il est réduit à
$\mathbf{F}_{2}$ en degré zéro. Donc il corespond à un schéma en groupes affine
$G$ sur $\mathbf{F}_{2}$, qui est d'ailleurs le foncteur des
$\otimes$-automorphismes additifs du foncteur $H^{\bullet}$ sur $C$ \add{qui
opèrent trivialement sur \struck{\ill{}} \uncertain{$H^{1}(S^{1})$}}.
\struck{\ill{}}\note{« i.e.\ commutative (car.\ 2 !) », « graduée
anticommutative donc » et « additifs » sont tapés en interligne et mis en place
par des traits à l'encre. L'exposant de $H(S^{1})$ est récrit à l'encre et se
lit $1$ sans certitude}

\page{48}

Le théorème de Thom est alors qu'un $x \in H^{\bullet}(X)$ est dans $S(X)$ sss
il est fixe par $G$, i.e.\ satisfait $A^{+} . x = 0$. (Pourquoi cette condition
est-elle déjà nécessaire ?) Une autre façon de le voir est de dire que le
foncteur $H^{\bullet}$ considéré comme allant de la catégorie $\mathbf{S}$
précédente dans la catégorie des représentations de $G$ dans des vectoriels sur
$\mathbf{F}_{2}$, est pleinement fidèle. Quand on prend l'enveloppe
pseudo-abélienne, on trouve une sous-catégorie pleine de la catégorie des
représentations de $G$, mais il est douteux qu'on trouve une catégorie
abélienne, voire la catégorie de toutes les représentations de $G$ de dimension
finie sur $\mathbf{F}_{2}$. (Quelle est la structure de $G$, qui est prolisse ;
d'après Quillen, il serait pro-unipotent ….). Ici $G$ bien sûr joue le rôle
d'un groupe de Galois motivique, mais contrairement à la situation familière, il
semble qu'il soit fortement non réductif (et d'ailleurs le corps de définition
dans tout ceci est $\mathbf{F}_{2}$, et non $\mathbf{Q}$ !).

\marginal{oui, car son anneau affine est une algèbre de polynômes à une
infinité d'indéterminées, \struck{\ill{}} réunion d'algèbres de polynômes à
\ill{} \ill{} \uncertain{finis} stables par $\Delta$}\note{à l'encre, de biais
dans la marge gauche, relié par un trait à « Quillen, il serait pro-unipotent »}

Quillent donne une généralisation de la construction d'une catégorie
universelle $\mathbf{B}$ au cas où on travaille avec des variétés pas
néessairement propres (mais $f_{\cdot}$ n'étant supposé défini que pour $f$
propre), et quand on introduit une condition d'orientation stable sur les
\emph{morphismes} envisagés, via un morphisme $H \to BO(\infty)$, où $H$ est un
H-espace (commutatif, associatif et unitaire ?) : une orientation de $f : X \to
Y$ sera un relèvement homotopique de $X \to BO$ défini par $T_{X} -
f^{\bullet}(T_{Y})$ en $X \to H$. les éléments de $H(X, Y)$ seront définis par
des diagrammes de morphismes orientés $Z \to X$, $Z \to Y$, le deuxième étant
propre, modulo bords de tels. En tant qu'ensemble, il se définit aussi en termes
de la variété $X \times Y$ munie simplement de la famille des supports formée
des ensembles fermés propres sur $Y$ …

\marginal{Les $B^{H}_{m}(X)$ ($X$ compact \uncertain{disons}) peuvent se décrire
homotopiquement à la Thom, comme les groupes d'homotopie stables d'espaces
$\underline{\mathrm{Hom}}(X, MH(N - m))$ - - - - -.}\note{à l'encre, à la suite de
la frappe, au pied de la page. Fin de la partie III et du tapuscrit dans ce
lot}

\page{58}

\marginal{$\dim E$}\note{le feuillet est la page 2 d'un tapuscrit en anglais de
D. G. Quillen, « On a conjecture of Adams », reproduit à l'alcool (pages 57 à
60), qui n'est pas transcrit. Seule y figure, d'une encre plus foncée que la
reproduction, cette addition, écrite en exposant du $p$ de la phrase « which is
of (inseparable) degree $p$ on each fiber », à propos de l'application
canonique $F : E \to E^{(p)}$ : le degré devient $p^{\dim E}$}

\end{document}
